\section{Non-Uniform Binary Fuse Filters}

This project was completed for CSC 592: Algorithms for Big Data.

\subsection{Motivation}

Distributed systems often require fast and small set data structures. For example, content delivery networks avoid caching one-hit wonders by querying a visited
set~\cite{maggs2015}. Assuming $n \ll |U|$, such exact membership queries take $\Omega(\lg |U|)$ bits per key, which is impractical for large $U$. By 
permitting false positive results with some probability $\epsilon$, the lower bound is reduced to $\lg \frac{1}{\epsilon}$ bits per key~\cite{carter1978}. That 
is, \textit{approximate} membership queries can be much more memory efficient than exact membership queries. \\

\noindent
A binary fuse filter~\cite{graf2022} is a recent probabilistic data structure for approximate membership queries. Membership information is stored in an array 
of $k$-bit values partitioned into disjoint segments. To query the filter, a key is hashed to a fingerprint value and three array entries located in consecutive 
segments (see Figure~\ref{csc592_bigdata:binaryfusefilter}). If the bitwise-XOR of the entries is equal to the fingerprint, then the query \emph{might} be in 
the set. Otherwise, the query \emph{definitely} is not in the set.

\begin{figure}[H]
    \centering
    \begin{tikzpicture}[>=Stealth, arraynode/.style={draw, minimum width=7mm, minimum height=6mm, outer sep=0, inner sep=0}]

        \node[arraynode, fill=lightyellow] (a1) {000};
        \node[arraynode, fill=lightorange, right=0 of a1] (a2) {011};
        \node[arraynode, fill=lightorange, right=0 of a2] (a3) {100};
        \node[arraynode, fill=lightyellow, right=0 of a3] (a4) {000};
        \node[arraynode, fill=lightyellow, right=0 of a4] (a5) {000};
        \node[arraynode, right=0 of a5, path picture={
        \shade[left color=lightorange, right color=lightred]
                (path picture bounding box.south west)
                rectangle (path picture bounding box.north east);
        }] (a6) {011};
        \node[arraynode, right=0 of a6, path picture={
        \shade[left color=lightred, right color=lightpink]
                (path picture bounding box.south west)
                rectangle (path picture bounding box.north east);
        }] (a7) {001};
        \node[arraynode, fill=lightyellow, right=0 of a7] (a8) {000};
        \node[arraynode, fill=lightyellow, right=0 of a8, path picture={
        \shade[left color=lightred, right color=lightpink]
                (path picture bounding box.south west)
                rectangle (path picture bounding box.north east);
        }] (a9) {001};
        \node[arraynode, fill=lightyellow, right=0 of a9] (a10) {000};
        \node[arraynode, fill=lightyellow, right=0 of a10, path picture={
        \shade[left color=lightpink, right color=lightpurple]
                (path picture bounding box.south west)
                rectangle (path picture bounding box.north east);
        }] (a11) {110};
        \node[arraynode, fill=lightyellow, right=0 of a11] (a12) {000};
        \node[arraynode, fill=lightpurple, right=0 of a12] (a13) {110};
        \node[arraynode, fill=lightyellow, right=0 of a13] (a14) {000};
        \node[arraynode, fill=lightyellow, right=0 of a14] (a15) {000};
        \node[arraynode, fill=lightpurple, right=0 of a15] (a16) {111};

        \draw[thick] 
            ([yshift=4pt]a1.north west) -- 
            ([yshift=-4pt]a1.south west);
        
        \foreach \i in {2,4,6,8,10,12,14,16} {
            \draw[thick] 
                ([yshift=4pt]a\i.north east) -- 
                ([yshift=-4pt]a\i.south east);
        }

        \node[circle, draw, above=12mm of a4, fill=lightorange] (x) {$x$};
        \draw[->] (x) to (a2.north);
        \draw[->] (x) to (a3.north);
        \draw[->] (x) to (a6.north);
        
        \node[circle, draw, above=12mm of a7, fill=lightred] (y) {$y$};
        \draw[->] (y) to (a6.north);
        \draw[->] (y) to (a7.north);
        \draw[->] (y) to (a9.north);

        \node[circle, draw, above=12mm of a10, fill=lightpink] (w) {$w$};
        \draw[->] (w) to (a7.north);
        \draw[->] (w) to (a9.north);
        \draw[->] (w) to (a11.north);
        
        \node[circle, draw, above=12mm of a13, fill=lightpurple] (z) {$z$};
        \draw[->] (z) to (a11.north);
        \draw[->] (z) to (a13.north);
        \draw[->] (z) to (a16.north);
        
    \end{tikzpicture}
    \caption{Binary fuse filter with 16-entry array, 2-entry segments, and 3-bit fingerprints. Vertical bars mark segment boundaries. The filter contains keys 
    $\mathcolorbox{lightorange}{x}$, $\mathcolorbox{lightred}{y}$, $\mathcolorbox{lightpink}{w}$ and $\mathcolorbox{lightpurple}{z}$, with fingerprints 
    $\mathcolorbox{lightorange}{100}$, $\mathcolorbox{lightred}{011}$, $\mathcolorbox{lightpink}{110}$ and $\mathcolorbox{lightpurple}{111}$, respectively. 
    Entries hashed to by keys are marked with arrows and the keys' colors. Each key hashes to three consecutive segments.}
    \label{csc592_bigdata:binaryfusefilter}
\end{figure}


\begin{figure}
\centering
\begin{tikzpicture}[scale=1.2]

\node (v1) at (0, 2) {};
\node (v2) at (0, 1) {};
\node (v3) at (2, 1.5) {};
\node (v4) at (4, 2) {};
\node (v5) at (4, 1) {};
\node (v6) at (6, 1.5) {};
\node (v7) at (8, 2) {};
\node (v8) at (8, 1) {};

\begin{scope}[fill opacity = 0.7]

\filldraw [fill=lightorange] 
    ($(v1) + (-1,0)$)
    to [out=90, in=180] ($(v1) + (0,1)$)
    to [out=0, in=90] ($(v3) + (1,0)$)
    to [out=270, in=0] ($(v2) + (0,-1)$)
    to [out=180, in=270] ($(v1) + (-1,0)$);

\filldraw [fill=lightred] 
    ($(v4) + (1,0)$)
    to [out=90, in=0] ($(v4) + (0,1)$)
    to [out=180, in=90] ($(v3) + (-1,0)$)
    to [out=270, in=180] ($(v5) + (0,-1)$)
    to [out=0, in=270] ($(v4) + (1,0)$);

\filldraw [fill=lightpink] 
    ($(v4) + (-1,0)$)
    to [out=90, in=180] ($(v4) + (0,1)$)
    to [out=0, in=90] ($(v6) + (1,0)$)
    to [out=270, in=0] ($(v5) + (0,-1)$)
    to [out=180, in=270] ($(v4) + (-1,0)$);

\filldraw [fill=lightpurple] 
    ($(v7) + (1,0)$)
    to [out=90, in=0] ($(v7) + (0,1)$)
    to [out=180, in=90] ($(v6) + (-1,0)$)
    to [out=270, in=180] ($(v8) + (0,-1)$)
    to [out=0, in=270] ($(v7) + (1,0)$);

\end{scope}

\foreach \i in {1,...,8}
    \fill (v\i) circle (0.1);

\foreach \i/\name in {1/A[1],2/A[2],3/A[5],4/A[6],5/A[8],6/A[10],7/A[12],8/A[15]}
    \node[right] at (v\i) {$\name$};

\end{tikzpicture}

\caption{Hypergraph induced by hash functions and keys from Figure~\ref{csc592_bigdata:binaryfusefilter}. Yellow hyperedge is induced by $x$, blue hyperedge
is induced by $y$, green hyperedge is induced by $w$, and red hyperedge is induced by $z$. Vertices with degree zero not shown.}

\end{figure}


\begin{figure}
\centering  
\[
\setlength{\arraycolsep}{2pt}
\begin{array}{cccccccccccccc}

\mathcolorbox{lightorange}{A[1]} 
&\oplus& \mathcolorbox{lightorange}{A[2]} 
&\oplus& \gradientmathbox{lightorange}{lightred}{A[5]} 
&=& \mathcolorbox{lightorange}{100} &\qquad

\gradientmathbox{lightorange}{lightred}{A[5]} 
&\oplus& \gradientmathbox{lightred}{lightpink}{A[6]} 
&\oplus& \gradientmathbox{lightred}{lightpink}{A[8]} 
&=& \mathcolorbox{lightred}{011}\\[1em]

\gradientmathbox{lightred}{lightpink}{A[6]} 
&\oplus& \gradientmathbox{lightred}{lightpink}{A[8]} 
&\oplus& \gradientmathbox{lightpink}{lightpurple}{A[10]} 
&=& \mathcolorbox{lightpink}{110} &\qquad

\gradientmathbox{lightpink}{lightpurple}{A[10]} 
&\oplus& \mathcolorbox{lightpurple}{A[12]} 
&\oplus& \mathcolorbox{lightpurple}{A[15]} 
&=& \mathcolorbox{lightpurple}{111}
\end{array}
\]

\caption{System of equations.}

\end{figure}

% TODO: binary fuse filter construction
% find assignment of values to array such that queries are correct for all elements in the filter.

% induces a system of linear equations in a finite field
% diagram?

% vertex == array location
% edge == key hashed to array locations

% induces random hypergraph family
% 2-peelablity: remove any vertex with degree 0 or 1, as well as incident edges
% peelable if iterating produces empty hypergraph
% phase transition, 2-peelablity threshold
% greedy algorithm

% spatial coupling ->
% bff construction is optimal for k-uniform families of hypergraphs (in limit of # of keys)
% no theorical improvements possible 



% 9% theory overhead, 13% practical overhead
% Optimality
Binary fuse filters 


\subsection{Research Questions}

% optimality of binary fuse for uniform, results with non-uniform

% TODO: revise for clarity, maybe parallel structure with filters not hypergraphs
\noindent
While fuse graph families saturate the uniform hypergraph peelability threshold, we are not aware of any constructions that saturate the 
peelability threshold for non-uniform hypergraph families. Therefore, we ask: 

\begin{enumerate}
    \item How can we systematically construct non-uniform hypergraph families with peelablity thresholds comparable to or surpassing those of fuse graph families?
    \item How can we design efficient approximate membership query data structures by exploiting peelable non-uniform hypergraph families?
\end{enumerate}

\subsection{Summary}


% correspondance between random hypergraph families and fuse filters
% binary fuse filters use uniform hypergraph families
% we considered non-uniform
% empirical peelablity thresholds for 10403 non-uniform fuse graph families
% mean edge size between 3 and 4

% can interpolate via sharding to attain any threshold

% non-uniform hypergraph families
% pareto frontier
% optimized implementation

% limitations:
% explore additional parameteratizations of hg families
% need to prove conjecture
% implementation did not perform well

TODO: A 400–500-word summary of the technical contributions of the work, what was
accomplished overall, and the limitations of their work.

\subsection{Reflection}

TODO: A 250-word self reflection of what skills were developed or enhanced in completing those
projects.